\section{Argument derivatives and all anchored primitive orders}
\label{clc:sec:primitives}

Order differentiation is distinct from differentiating the center or
the logarithmic argument. The contour formula makes the latter two
operations equally explicit:
\begin{align}
 \partial_A\CC_{\bq,A}(W;\mu)
      &=-W\CC_{\bq,A}(W+1;\mu),\label{clc:eq:A-ladder}\\
 (\partial_\mu+A)\CC_{\bq,A}(W;\mu)
      &=\CC_{\bq,A}(W-1;\mu).
 \label{clc:eq:mu-ladder}
\end{align}
Both follow by differentiating \eqref{clc:eq:C-contour}; its exponential
vertical majorant justifies each derivative. They also follow from the
corresponding profile identities. Iteration gives
\begin{equation}\label{clc:eq:A-higher}
 \partial_A^d\CC_{\bq,A}(W;\mu)
       =(-1)^d(W)_d\CC_{\bq,A}(W+d;\mu).
\end{equation}
Every mixed spectral derivative can be applied to these identities.
The Pochhammer factor must then be differentiated too.

\subsection{First and higher anchored antiderivatives}
For an anchor $A_0$ in $\Rea A>-q_*$,
\begin{equation}\label{clc:eq:first-primitive}
 \int_{A_0}^A\CC_{\bq,u}(W;\mu)\dd u
 =\frac{\CC_{\bq,A_0}(W-1;\mu)
                  -\CC_{\bq,A}(W-1;\mu)}{W-1}.
\end{equation}
The integral is along any path in that half-plane. The right side is
filled in at $W=1$ by
\begin{equation}\label{clc:eq:first-resonant}
 \CC_{\bq,A_0}^{[1]}(0;\mu)-\CC_{\bq,A}^{[1]}(0;\mu),
\end{equation}
since $\CC_{\bq,A}(0;\mu)=K_{\bq}(\mu)$ is independent of $A$.
At $\mu=0$, this already gives a finite Gamma and polygamma primitive
of the weight-one convolution.

The same mechanism supplies every primitive order. Define
\begin{align}
 T_d(W;A,A_0;\mu)={}&\CC_{\bq,A}(W-d;\mu)\notag\\
 &-\sum_{j=0}^{d-1}\frac{(A-A_0)^j}{j!}(-1)^j(W-d)_j
             \CC_{\bq,A_0}(W-d+j;\mu).
 \label{clc:eq:T-d}
\end{align}

\begin{theorem}[All anchored primitives, including resonance]
\label{clc:thm:primitives}
For $d\ge1$,
\begin{equation}\label{clc:eq:primitive-d}
 \boxed{\displaystyle
 \int_{A_0}^A\frac{(A-u)^{d-1}}{(d-1)!}
                 \CC_{\bq,u}(W;\mu)\dd u
 =\frac{(-1)^dT_d(W;A,A_0;\mu)}{(W-d)_d}.}
\end{equation}
The right side is entire in $W$ after removing its apparent poles at
$W=1,\ldots,d$. At $W=k$ in that range its value is
\begin{equation}\label{clc:eq:primitive-resonance}
 \frac{(-1)^k}{(k-1)!(d-k)!}
            \left.\partial_WT_d(W;A,A_0;\mu)\right|_{W=k}.
\end{equation}
All order derivatives of these primitive identities are valid.
\end{theorem}
\begin{proof}
Apply the ordinary Taylor formula with integral remainder, in the
variable $A$, to $\CC_{\bq,A}(W-d;\mu)$ at $A_0$. Its $j$th derivative
is given by \eqref{clc:eq:A-higher}. Its $d$th derivative is
$(-1)^d(W-d)_d\CC_{\bq,A}(W;\mu)$. This proves
\eqref{clc:eq:primitive-d} when the polynomial denominator is nonzero.
The left side is entire in $W$, by local uniform holomorphy on a compact
integration path. Therefore all apparent poles are removable.
The polynomial $(W-d)_d=\prod_{j=1}^d(W-j)$ has simple zeros, with
its derivative at $W=k$ equal to
$(-1)^{d-k}(k-1)!(d-k)!$.
Taking the quotient limit gives \eqref{clc:eq:primitive-resonance}.
Local uniform holomorphy and Cauchy's formula give every order derivative.
\end{proof}

This formulation fixes all integration constants by the anchor and
its vanishing derivatives through order $d-1$. It is stronger than
writing an unanchored formal primitive and leaving its constants
ambiguous at resonant orders.

\subsection{Explicit Gamma primitives of the two-scale family}
For $D_m(W;A)=\CC_{(1,m),A}(W;0)$, define $b_r=(A+r)/m$.
When $m$ is odd, put
\begin{equation}\label{clc:eq:Podd}
 P_m(A)=\frac1m\left[\psi(b_m)
                  -\pi\sum_{r=1}^{m-1}c_r^{(m)}\log\Gamma(b_r)\right].
\end{equation}
When $m$ is even, put
\begin{equation}\label{clc:eq:Peven}
 P_m(A)=\frac1m\left[\eta(1,b_m)-\pi\sum_{r=1}^{m-1}c_r^{(m)}
 \log\frac{\Gamma(b_r/2)}{\Gamma((b_r+1)/2)}\right].
\end{equation}
Then in both cases
\begin{equation}\label{clc:eq:Pderivative}
 P_m'(A)=D_m(1;A),\qquad
 \int_{A_0}^AD_m(1;u)\dd u=P_m(A)-P_m(A_0).
\end{equation}
To check the odd case directly, use
$\psi'(b)=\zeta(2,b)$ and
$\sum_rc_r^{(m)}\zeta(1,b_r)=-\sum_rc_r^{(m)}\psi(b_r)$,
where the sum is understood after its residue cancellation.
For the even case use
$\partial_b\eta(1,b)=-\eta(2,b)$ and
\[
 \partial_b\log\frac{\Gamma(b/2)}{\Gamma((b+1)/2)}=-\eta(1,b).
\]
Both derivatives reproduce the $W=1$ case of the normalized pair formula,
including its $m^{-2}$ factor. This is also an independent check of the
sign in \eqref{clc:eq:first-resonant}.

Subsequent center derivatives of \eqref{clc:eq:Podd} and
\eqref{clc:eq:Peven} are finite polygamma expressions. The all-index
Stieltjes formulas in Section~\ref{clc:sec:stieltjes}, followed by
\eqref{clc:eq:primitive-d}, give their anchored order-jet primitives.
Thus differentiation and integration are both covered without replacing
a pole-cancelled expression by its singular separate pieces.
